\documentclass[a4paper, 14pt]{extarticle}

% Поля
%--------------------------------------
\usepackage{geometry}
\geometry{a4paper,tmargin=2cm,bmargin=2cm,lmargin=3cm,rmargin=1cm}
%--------------------------------------

\usepackage{float}
%Russian-specific packages
%--------------------------------------
\usepackage[T2A]{fontenc}
\usepackage[utf8]{inputenc} 
\usepackage[english, main=russian]{babel}
%--------------------------------------

\usepackage{textcomp}
\usepackage{xcolor}

% Красная строка
%--------------------------------------
\usepackage{indentfirst}               
%--------------------------------------
%Листинги
%--------------------------------------
\usepackage{listings}
\lstset{
  basicstyle=\ttfamily\footnotesize,
  breakatwhitespace=false,
  breaklines=true,
  captionpos=t,
  inputencoding=utf8,
  extendedchars=true,
  literate={а}{а}1{б}{б}1{в}{в}1{г}{г}1{д}{д}1{е}{е}1{ё}{ё}1{ж}{ж}1{з}{з}1{и}{и}1{й}{й}1{к}{к}1{л}{л}1{м}{м}1{н}{н}1{о}{о}1{п}{п}1{р}{р}1{с}{с}1{т}{т}1{у}{у}1{ф}{ф}1{х}{х}1{ц}{ц}1{ч}{ч}1{ш}{ш}1{щ}{щ}1{ъ}{ъ}1{ы}{ы}1{ь}{ь}1{э}{э}1{ю}{ю}1{я}{я}1,
  frame=single,
  keepspaces=true,
  keywordstyle=\bf,
  numbers=none,
  xleftmargin=25pt,
  xrightmargin=25pt,
  showspaces=false,
  showstringspaces=false,
  showtabs=false,
  stepnumber=1,
  tabsize=2,
  title=\lstname
}
%--------------------------------------
\usepackage[most]{tcolorbox}
\tcbuselibrary{listings,breakable}
\newtcblisting{breakablelisting}{
  listing only,
  breakable,
  enhanced,
  frame hidden,
  colback=white,
  left=8pt,
  overlay={\draw[line width=0.4pt] (frame.north west) -- (frame.south west);
          \draw[line width=0.4pt] ([xshift=3pt]frame.north west) -- ([xshift=3pt]frame.south west);},
  listing options={
    frame=none,
    numbers=none,
    mathescape=false,
    texcl=false,
    inputencoding=utf8,
    extendedchars=true,
    literate={а}{а}1{б}{б}1{в}{в}1{г}{г}1{д}{д}1{е}{е}1{ё}{ё}1{ж}{ж}1{з}{з}1{и}{и}1{й}{й}1{к}{к}1{л}{л}1{м}{м}1{н}{н}1{о}{о}1{п}{п}1{р}{р}1{с}{с}1{т}{т}1{у}{у}1{ф}{ф}1{х}{х}1{ц}{ц}1{ч}{ч}1{ш}{ш}1{щ}{щ}1{ъ}{ъ}1{ы}{ы}1{ь}{ь}1{э}{э}1{ю}{ю}1{я}{я}1,
    literate={\$}{{\$}}1
  }
}


%Graphics
%--------------------------------------
\usepackage{graphicx}
\usepackage{wrapfig}
%--------------------------------------
\graphicspath{{./lab5/}{./}}
% Полуторный интервал
%--------------------------------------
\linespread{1.3}                    
%--------------------------------------

%Выравнивание и переносы
%--------------------------------------
% Избавляемся от переполнений
\sloppy
% Запрещаем разрыв страницы после первой строки абзаца
\clubpenalty=10000
% Запрещаем разрыв страницы после последней строки абзаца
\widowpenalty=10000
%--------------------------------------

%Списки
\usepackage{enumitem}

%Подписи
\usepackage{caption} 

%Гиперссылки
\usepackage{hyperref}

\hypersetup {
	unicode=true
}

%Рисунки
%--------------------------------------
\DeclareCaptionLabelSeparator*{emdash}{~--- }
\captionsetup[figure]{labelsep=emdash,font=onehalfspacing,position=bottom}
%--------------------------------------

\usepackage{tempora}

%--------------------------------------

%%% Математические пакеты %%%
%--------------------------------------
\usepackage{amsthm,amsfonts,amsmath,amssymb,amscd}  % Математические дополнения от AMS
\usepackage{mathtools}                              % Добавляет окружение multlined
\usepackage[perpage]{footmisc}
%--------------------------------------

%--------------------------------------
%			НАЧАЛО ДОКУМЕНТА
%--------------------------------------

\begin{document}

%--------------------------------------
%			ТИТУЛЬНЫЙ ЛИСТ
%--------------------------------------
\begin{titlepage}
\thispagestyle{empty}
\newpage


%Шапка титульного листа
%--------------------------------------
\vspace*{-60pt}
\hspace{-65pt}
\begin{minipage}{0.3\textwidth}
\hspace*{-20pt}\centering
\includegraphics[width=\textwidth]{emblem}
\end{minipage}
\begin{minipage}{0.67\textwidth}\small \textbf{
\vspace*{-0.7ex}
\hspace*{-6pt}\centerline{Министерство науки и высшего образования Российской Федерации}
\vspace*{-0.7ex}
\centerline{Федеральное государственное автономное образовательное учреждение }
\vspace*{-0.7ex}
\centerline{высшего образования}
\vspace*{-0.7ex}
\centerline{<<Московский государственный технический университет}
\vspace*{-0.7ex}
\centerline{имени Н.Э. Баумана}
\vspace*{-0.7ex}
\centerline{(национальный исследовательский университет)>>}
\vspace*{-0.7ex}
\centerline{(МГТУ им. Н.Э. Баумана)}}
\end{minipage}
%--------------------------------------

%Полосы
%--------------------------------------
\vspace{-25pt}
\hspace{-35pt}\rule{\textwidth}{2.3pt}

\vspace*{-20.3pt}
\hspace{-35pt}\rule{\textwidth}{0.4pt}
%--------------------------------------

\vspace{1.5ex}
\hspace{-35pt} \noindent \small ФАКУЛЬТЕТ\hspace{80pt} <<Информатика и системы управления>>

\vspace*{-16pt}
\hspace{47pt}\rule{0.83\textwidth}{0.4pt}

\vspace{0.5ex}
\hspace{-35pt} \noindent \small КАФЕДРА\hspace{50pt} <<Теоретическая информатика и компьютерные технологии>>

\vspace*{-16pt}
\hspace{30pt}\rule{0.866\textwidth}{0.4pt}
  
\vspace{11em}

\begin{center}
\Large {\bf Домашняя работа №1} \\ 
\large {\bf по курсу <<Моделирование>>} \\
\end{center}\normalsize

\vspace{8em}


\begin{flushright}
  {Студент группы ИУ9-82Б Шемякин В.А. \hspace*{15pt}\\ 
  \vspace{2ex}
  Преподаватель Домрачева А.Б.\hspace*{15pt}}
\end{flushright}

\bigskip

\vfill
 

\begin{center}
\textsl{Москва 2026}
\end{center}
\end{titlepage}
%--------------------------------------
%		КОНЕЦ ТИТУЛЬНОГО ЛИСТА
%--------------------------------------

\renewcommand{\ttdefault}{pcr}

\setlength{\tabcolsep}{3pt}
\newpage
\setcounter{page}{2}
\section{Постановка задачи}

Студенты первого курса сдают зачёт по информатике, имея три попытки. Вероятность сдать зачёт с первой попытки равна 0,6, со второй~--- 0,25, с третьей~--- 0,1. В один день зачёт могут сдавать не более 20 человек. Между двумя попытками должно пройти не менее одного дня. Требуется провести вычислительный эксперимент с целью оценки количества студентов, успешно сдавших зачёт за 6 рабочих дней зачётной недели. В эксперименте участвует 80 студентов.

\section{Построение модели}

Модель представляет собой систему массового обслуживания с многоканальным обслуживанием и приоритетной очередью. Заявками в модели являются студенты: каждая попытка сдачи экзамена рассматривается как одна заявка на обслуживание. Источник заявок~--- поток студентов, поступающих на сдачу зачёта. Очередь имеет ограниченную дисциплину: приоритет отдаётся пересдающим. Номер попытки $P_1 \in \{1, 2, 3\}$ задаёт приоритет по формуле $\text{priority} = 4 - P_1$: чем выше номер попытки, тем выше приоритет.

Каналы обслуживания~--- 20 мест для сдачи экзамена: одновременно экзамен могут сдавать не более 20 человек. Время обслуживания~--- 1440 минут (один рабочий день). Заявка, не прошедшая обслуживание (несдача), после задержки 1440 минут возвращается в очередь с увеличенным номером попытки.

Компоненты системы: источник заявок, очередь с приоритетом, устройство обслуживания. Связи: заявки из источника поступают в очередь, при наличии свободного канала занимают его, проходят обслуживание, по результату (успех или неудача) либо покидают систему, либо после задержки возвращаются в очередь.

Вероятности успеха: $p_1 = 0{,}6$, $p_2 = 0{,}25$, $p_3 = 0{,}1$. Ожидаемое число сдавших при неограниченной ёмкости:
\begin{equation}
M = n \cdot \bigl( p_1 + (1-p_1) p_2 + (1-p_1)(1-p_2) p_3 \bigr),
\end{equation}
где $n$~--- число студентов. При $n = 80$ теоретическое ожидание $M \approx 58$.

\section{Решение вычислительных задач}

Поток заявок формируется источником с интервалом между появлениями, заданным равномерным распределением на отрезке $[48, 96]$ минут и лимитом 80 заявок. Таким образом, 80 студентов поступают в систему постепенно в начале зачётной недели.

Работа потоков заявок: первичный поток направляется в очередь на экзамен. Каждая заявка при входе получает приоритет по номеру попытки и ожидает свободный канал. После обслуживания заявка либо завершается (сдача), либо через задержку 1440 минут возвращается в очередь. Поток пересдач образует возвратные заявки с приоритетом выше первичных. Дополнительно в модели присутствует поток-ограничитель: таймер создаётся в начальный момент, выдерживает 10080 минут и принудительно завершает симуляцию.

\section{Анализ модели}

\refstepcounter{lstlisting}\label{lst:output}
Листинг~\ref{lst:output}. Вывод работы программы.

\begin{lstlisting}
              GPSS World Simulation Report - Untitled Model 1.3.1

                   Monday, March 09, 2026 21:15:08  

           START TIME           END TIME  BLOCKS  FACILITIES  STORAGES
                0.000          10080.000    27        0          1

              NAME                       VALUE  
          ATTEMPT                         3.000
          ATTEMPT2                       11.000
          ATTEMPT3                       13.000
          EXHAUSTED                      22.000
          FAILED                         17.000
          NEXT_DAY                       19.000
          PASS                        10007.000
          STOR                        10000.000
          SUCCESS                        14.000
          TOO_LATE                       23.000

 LABEL              LOC  BLOCK TYPE     ENTRY COUNT CURRENT COUNT RETRY
                    1    GENERATE            80             0       0
                    2    ASSIGN              80             0       0
ATTEMPT             3    PRIORITY           127             0       0
                    4    QUEUE              127             0       0
                    5    ENTER              127             0       0
                    6    TEST               127             0       0
                    7    DEPART             111             0       0
                    8    ADVANCE            111             0       0
                    9    TEST               111             0       0
                   10    TRANSFER            80             0       0
ATTEMPT2           11    TEST                31             0       0
                   12    TRANSFER            30             0       0
ATTEMPT3           13    TRANSFER             1             0       0
SUCCESS            14    SAVEVALUE           52             0       0
                   15    LEAVE               52             0       0
                   16    TERMINATE           52             0       0
FAILED             17    LEAVE               59             0       0
                   18    TEST                59             0       0
NEXT_DAY           19    ADVANCE             58            11       0
                   20    ASSIGN              47             0       0
                   21    TRANSFER            47             0       0
EXHAUSTED          22    TERMINATE            1             0       0
TOO_LATE           23    LEAVE               16             0       0
                   24    TERMINATE           16             0       0

QUEUE              MAX CONT. ENTRY ENTRY(0) AVE.CONT. AVE.TIME   AVE.(-0) RETRY
 1                  18   16    127     43     5.659    449.130    679.042   0

STORAGE            CAP. REM. MIN. MAX.  ENTRIES AVL.  AVE.C. UTIL. RETRY DELAY
 STOR               20   20   0    20      127   1   15.857  0.793    0    0

SAVEVALUE               RETRY       VALUE
 PASS                     0         52.000
\end{lstlisting}

За 6 рабочих дней (8640 минут) и один дополнительный день (до 10080 минут) для завершения экзаменов, начатых в конце недели, модель обработала 80 студентов. Сдали зачёт 52 человека. Всего экзаменационных попыток 111 из максимума 120. 16 студентов получили слот на экзамен после окончания зачётной недели и не были допущены к сдаче. 59 студентов не сдали с первой или второй попытки и перешли к ожиданию следующего дня. 11 заявок на момент окончания симуляции ожидали пересдачу. 1 студент исчерпал три попытки. Загрузка каналов составила 79,3\%.

\section{Вывод}
В ходе работы была построена модель системы массового обслуживания. Модель была реализована в среде GPSS, а также был проведен анализ модели. Также был проведен вычислительный эксперимент для 80 студентов. Число сдавших зачет оказалось равным 52.

\section{Практическая реализация}

\refstepcounter{lstlisting}\label{lst:gpss}
Листинг~\ref{lst:gpss}. Исходный код модели.

\begin{lstlisting}
SIMULATE

STOR    STORAGE 20

        GENERATE 72,24,,80
        ASSIGN 1,1

ATTEMPT PRIORITY (4-P1)
        QUEUE 1
        ENTER STOR,1
        TEST L AC1,8640,TOO_LATE
        DEPART 1
        ADVANCE 1440
        TEST E P1,1,ATTEMPT2
        TRANSFER 0.6,FAILED,SUCCESS

ATTEMPT2 TEST E P1,2,ATTEMPT3
        TRANSFER 0.25,FAILED,SUCCESS

ATTEMPT3 TRANSFER 0.1,FAILED,SUCCESS

SUCCESS SAVEVALUE PASS+,1
        LEAVE STOR,1
        TERMINATE 1

FAILED  LEAVE STOR,1
        TEST L P1,3,EXHAUSTED
NEXT_DAY ADVANCE 1440
        ASSIGN 1+,1
        TRANSFER ,ATTEMPT

EXHAUSTED TERMINATE 1

TOO_LATE LEAVE STOR,1
        TERMINATE 1

        GENERATE ,,0,1
        ADVANCE 10080
        TERMINATE 1000

        RMULT 12345
        START 80
\end{lstlisting}
\end{document}
